% Part: incompleteness
% Chapter: introduction
% Section: definitions

\documentclass[../../../include/open-logic-section]{subfiles}

\begin{document}

\olfileid[fr]{inc}{int}{def}

\olsection{Définitions}

Pour mener à bien le projet de Hilbert, qui consistait à formaliser
les mathématiques puis à montrer que cette formalisation est cohérente
et complète, il faut d'abord choisir un langage, un cadre logique et
un système d'axiomes. Supposons, pour notre propos, que les
mathématiques puissent être formalisées dans un langage du premier
ordre : il existe alors un ensemble de !!{constant}s, de
!!{function}s et de !!{predicate}s qui, avec les connecteurs et les
quantificateurs de la logique du premier ordre, permettent d'exprimer
les assertions mathématiques. On admet généralement l'existence d'un
tel langage : celui de la théorie des ensembles, dont~$\in$ est le
seul symbole non logique. Qu'un langage si simple soit aussi expressif
semble d'abord peu plausible; il a fallu beaucoup de travail pour
établir que presque toutes les mathématiques peuvent se formuler dans
ce vocabulaire austère. Pour simplifier, limitons-nous à présent à
l'arithmétique, c'est-à-dire à la partie des mathématiques qui traite
des nombres naturels~$\Nat$. Nous retiendrons pour exprimer les
faits arithmétiques le langage~$\Lang L_A$. Il contient un seul !!{predicate}
binaire~$<$, une seule !!{constant}~$\Obj 0$, un !!{function} unaire
$\prime$, et deux !!{function}s binaires~$+$ et~$\times$.

\begin{defn}
Un ensemble d'!!{sentence}s~$\Gamma$ est une \emph{théorie} s'il est
fermé par conséquence logique, c'est-à-dire si
$\Gamma=\Setabs{!A}{\Gamma\Entails !A}$.
\end{defn}

On peut définir une théorie de deux façons simples. La première
consiste à prendre l'ensemble des !!{sentence}s vrais dans une
!!{structure}. Considérons, par exemple, la !!{structure} du
langage~$\Lang L_A$ dont le !!{domain} est~$\Nat$ et dans laquelle
chaque symbole non logique reçoit son interprétation usuelle.

\begin{defn}\ollabel{def:standard-model}
Le \emph{modèle standard de l'arithmétique} est la !!{structure}~$\Struct
N$ définie ainsi :
\begin{enumerate}
\item $\Domain N = \Nat$
\item $\Assign{\Obj 0}{N} = 0$
\item $\Assign{\Obj \prime}{N}(n) = n + 1$ pour tout $n \in \Nat$
\item $\Assign{\Obj +}{N}(n, m) = n + m$ pour tous $n, m \in \Nat$
\item $\Assign{\Obj \times}{N}(n, m) = n\cdot m$ pour tous $n, m \in \Nat$
\item $\Assign{\Obj <}{N} = \Setabs{\tuple{n, m}}{n \in \Nat, m \in
  \Nat, n < m}$
\end{enumerate}
\end{defn}

Distinguons $\times$ de~$\cdot$. Le premier est un symbole du langage
arithmétique. Il évoque certes la multiplication, mais ce n'est pas
l'opération de multiplication : c'est un symbole de fonction binaire
(officiellement, $\Obj f^2_1$). En revanche, $\cdot$ est la fonction
ordinaire de multiplication. Ainsi, $n\cdot m$ désigne le produit des
nombres $n$ et~$m$, tandis que $x\times y$ est un terme du langage
arithmétique. Dans le modèle standard, le symbole de fonction~$\times$
est interprété par la fonction~$\cdot$ sur les nombres naturels.
Pour l'addition, nous employons~$+$ aussi bien pour le symbole de
fonction du langage que pour l'opération sur les nombres naturels;
le contexte indique lequel des deux est visé.

\begin{defn}
La théorie de l'\emph{arithmétique vraie} est l'ensemble des
!!{sentence}s satisfaits dans le modèle standard de l'arithmétique :
\[
\Th{TA} = \Setabs{!A}{\Sat{N}{!A}}.
\]
\end{defn}

$\Th{TA}$ est bien une théorie. En effet, si $\Th{TA}\Entails !A$,
alors~$!A$ est satisfait dans toute !!{structure} qui satisfait
$\Th{TA}$. Or $\Sat{N}{\Th{TA}}$; donc $\Sat{N}{!A}$ et
$!A\in\Th{TA}$.

La seconde façon de définir une théorie~$\Gamma$ consiste à prendre
l'ensemble des !!{sentence}s qui découlent d'un ensemble
d'énoncés~$\Gamma_0$. $\Gamma$ est alors la « clôture » de~$\Gamma_0$
par conséquence logique. Cette présentation n'est intéressante que
si~$\Gamma_0$ est donné explicitement, par exemple par une liste ou
un schéma d'axiomes. Pour les axiomatisations effectives étudiées ici, nous demandons que~$\Gamma_0$ soit
décidable\footnote{Précision éditoriale : la source présente cette exigence comme une condition de toute axiomatisation. La définition abstraite qui suit n'impose aucune effectivité à~$\Gamma_0$ ; celle-ci intervient dans la notion d'axiomatisabilité retenue plus bas.} : il doit exister un test calculable permettant de savoir
si un !!{sentence} appartient à~$\Gamma_0$. Les !!{sentence}s
de~$\Gamma_0$ sont appelés \emph{axiomes} de~$\Gamma$, et l'on dit
que~$\Gamma$ est \emph{axiomatisée} par~$\Gamma_0$.

\begin{defn}
Une théorie~$\Gamma$ est \emph{axiomatisée} par~$\Gamma_0$ si et
seulement si
\[
\Gamma = \Setabs{!A}{\Gamma_0 \Entails !A}
\]
\end{defn}

\begin{defn}
La théorie $\Th{Q}$ axiomatisée par les !!{sentence}s suivants
s'appelle \emph{arithmétique de Robinson}; c'est une théorie
arithmétique très simple.
\begin{align*}
& \lforall[x][\lforall[y][(\eq[x'][y'] \lif \eq[x][y])]] \tag{$!Q_1$}\\
& \lforall[x][\eq/[\Obj 0][x']] \tag{$!Q_2$}\\
& \lforall[x][(\eq[x][\Obj 0] \lor \lexists[y][\eq[x][y']])] \tag{$!Q_3$}\\
& \lforall[x][\eq[(x + \Obj 0)][x]] \tag{$!Q_4$}\\
& \lforall[x][\lforall[y][\eq[(x + y')][(x + y)']]] \tag{$!Q_5$}\\
& \lforall[x][\eq[(x \times \Obj 0)][\Obj 0]] \tag{$!Q_6$}\\
& \lforall[x][\lforall[y][\eq[(x \times y')][((x \times y) + x)]]] \tag{$!Q_7$}\\
& \lforall[x][\lforall[y][(x < y \liff \lexists[z][\eq[(z' + x)][y]])]] \tag{$!Q_8$}
\end{align*}
Les !!{sentence}s $\{!Q_1,\dots,!Q_8\}$ sont les axiomes de
$\Th{Q}$; cette théorie contient donc tous les !!{sentence}s qui
en découlent :
\[
\Th{Q} = \Setabs{!A}{\{!Q_1, \dots, !Q_8\} \Entails !A}.
\]
\end{defn}

\begin{defn}
Soit $!A(x)$ une !!{formula} de~$\Lang L_A$ dont les variables libres
sont~$x$ et $y_1$, \dots, $y_n$. Tout !!{sentence} de la forme
\[
\lforall[y_1][\dots\lforall[y_n][((!A(\Obj 0) \land \lforall[x][(!A(x)
\lif !A(x'))]) \lif \lforall[x][!A(x)])]]
\]
est une instance du \emph{schéma d'induction}. Les paramètres
$y_1$, \dots, $y_n$ de~$!A$ sont sous-entendus dans les occurrences
$!A(\Obj 0)$, $!A(x)$ et $!A(x')$.

L'\emph{arithmétique de Peano}~$\Th{PA}$ est la théorie axiomatisée
par les axiomes de~$\Th{Q}$ et toutes les instances du schéma
d'induction.
\end{defn}

\begin{explain}
Toute instance du schéma d'induction est vraie dans~$\Struct{N}$.
Le cas le plus simple est celui où la !!{formula}~$!A$ n'a que la
variable libre~$x$. Alors $!A(x)$ définit un sous-ensemble~$X_{!A}$
de~$\Nat$ dans~$\Struct{N}$ : il contient précisément les
$n\in\Nat$ tels que $\Sat{N}{!A(x)}[s]$ lorsque $s(x)=n$.
L'instance correspondante du schéma est
\[
((!A(\Obj 0) \land \lforall[x][(!A(x) \lif !A(x'))]) \lif
  \lforall[x][!A(x)]).
\]
Si son antécédent est vrai dans~$\Struct{N}$, alors $0\in X_{!A}$
et, dès que $n\in X_{!A}$, on a aussi $n+1\in X_{!A}$. De
$0\in X_{!A}$ on déduit $1\in X_{!A}$, puis $2\in X_{!A}$,
et ainsi de suite. Donc tout $n\in\Nat$ appartient à~$X_{!A}$,
ce qui signifie que $\lforall[x][!A(x)]$ est satisfait
dans~$\Struct{N}$.
\end{explain}

$\Th{Q}$ et $\Th{PA}$ sont toutes deux des théories axiomatisées.
Mais quelle est leur force ? $\Th{PA}$ peut-elle, par exemple,
démontrer toutes les vérités sur~$\Nat$ exprimables dans~$\Lang L_A$ ?
Ses axiomes tranchent-ils toutes les questions formulables dans ce
langage ?

Autrement dit, a-t-on $\Th{PA}=\Th{TA}$ ? La théorie~$\Th{TA}$
contient, par définition, toutes les vérités sur~$\Nat$ et tranche
toutes les questions exprimables dans~$\Lang L_A$ : pour tout
!!{sentence}~$!A$ de ce langage, on a soit $\Sat{N}{!A}$, soit
$\Sat{N}{\lnot !A}$, et donc soit $\Th{TA}\Entails !A$, soit
$\Th{TA}\Entails\lnot !A$. On appelle une telle théorie
\emph{complète}.

\begin{defn}
Une théorie~$\Gamma$ est \emph{complète} si et seulement si, pour
chaque !!{sentence}~$!A$ de son langage, on a
$\Gamma\Entails !A$ ou $\Gamma\Entails\lnot !A$.
\end{defn}

\begin{explain}
Par le théorème de complétude, $\Gamma\Entails !A$ si et seulement
si $\Gamma\Proves !A$. Ainsi, une théorie~$\Gamma$ est complète si
et seulement si, pour chaque !!{sentence}~$!A$ de son langage,
$\Gamma\Proves !A$ ou $\Gamma\Proves\lnot !A$.
\end{explain}

Une autre question se pose : existe-t-il une procédure calculable
qui permette de déterminer si un !!{sentence} appartient
à~$\Th{TA}$, à~$\Th{PA}$, ou même seulement à~$\Th{Q}$ ? Pour
la formuler précisément, définissons ce que signifie être
!!{decidable} pour un ensemble, par exemple un ensemble
d'!!{sentence}s.

\begin{defn}
Un ensemble~$X$ est \emph{!!{decidable}} s'il existe une procédure
calculable qui, recevant~$x$ en entrée, renvoie~$1$ si $x\in X$
et~$0$ sinon.
\end{defn}

La question devient donc : $\Th{TA}$, $\Th{PA}$ ou~$\Th{Q}$
sont-elles !!{decidable}s ?

La réponse est négative dans les trois cas. Mais ce phénomène ne
concerne pas seulement ces théories : \emph{toute} théorie qui
satisfait certaines conditions y est soumise. L'une d'elles,
vérifiée par $\Th{Q}$ et~$\Th{PA}$, est de posséder un ensemble
!!{decidable} d'axiomes.

\begin{defn}
Une théorie est \emph{!!{axiomatizable}} si elle admet un ensemble
!!{decidable} d'axiomes qui l'axiomatise.
\end{defn}

\begin{ex}
Toute théorie axiomatisée par un ensemble fini d'!!{sentence}s est
!!{axiomatizable}, car tout ensemble fini est !!{decidable}.
$\Th{Q}$, par exemple, est donc !!{axiomatizable}.

Les théories axiomatisées par schéma, comme~$\Th{PA}$, sont elles
aussi !!{axiomatizable}s. Pour tester si~$!B$ est un axiome
de~$\Th{PA}$, c'est-à-dire calculer la fonction~$\Char{X}$ qui vaut
$\Char{X}(!B)=1$ si~$!B$ en est un et~$0$ sinon, procédons ainsi.
Vérifions d'abord si~$!B$ est l'un des axiomes de~$\Th{Q}$ : si oui,
la réponse est « oui » et $\Char{X}(!B)=1$. Sinon, vérifions si~$!B$
est une instance du schéma d'induction. Cela se fait mécaniquement.
Toute instance commence par un nombre éventuellement nul de
quantificateurs universels, suivi d'une sous-!!{formula} qui est une
implication. Le conséquent de cette implication est
$\lforall[x][!A(x,y_1,\dots,y_n)]$, où $x,y_1,\dots,y_n$ sont les
variables libres de~$!A$; les quantificateurs initiaux de~$!B$
lient~$y_1,\dots,y_n$. Après avoir extrait~$!A$ et vérifié que ses
variables libres correspondent à celles liées par les quantificateurs
initiaux et par~$\lforall[x]$, on vérifie que l'antécédent est bien
\[
!A(\Obj 0, y_1, \dots, y_n) \land \lforall[x][(!A(x, y_1, \dots, y_n)
\lif !A(x', y_1, \dots, y_n))]
\]
Si c'est le cas, $!B$ est une instance du schéma; sinon,
il n'en est pas une.
\end{ex}

Pour répondre à cette question, et plus généralement pour savoir
quelles théories sont complètes ou décidables, la définition suivante
sera utile. Rappelons qu'un ensemble~$X$ est !!{enumerable} s'il est
vide ou s'il existe une fonction surjective~$f\colon\Nat\to X$;
on appelle cette fonction une \emph{énumération} de~$X$.

\begin{defn}
Un ensemble~$X$ est \emph{!!{computably enumerable}}
(\emph{e.é.} en abrégé) si et seulement s'il est vide ou admet une
énumération calculable.
\end{defn}

Outre l'axiomatisabilité, une autre condition nécessaire à
l'application des théorèmes d'incomplétude est que les théories
soient assez fortes pour démontrer des faits élémentaires sur les
fonctions calculables et les relations !!{decidable}s. Par « faits
élémentaires », nous entendons des !!{sentence}s qui expriment les
valeurs des fonctions calculables pour chacun de leurs arguments.
Être « assez forte » signifie que la théorie compte ces
!!{sentence}s parmi ses théorèmes. Prenons, par exemple, l'addition.
La valeur de~$+$ pour les arguments~$2$ et~$3$ est~$5$, soit
$2+3=5$. Le langage arithmétique exprime ce fait par l'!!{sentence}
$(\num{2}+\num{3})=\num{5}$, que~$\Th{Q}$ démontre.
Plus généralement, pour chaque fonction calculable~$f(x_1,x_2)$,
nous voudrions disposer d'une !!{formula}~$!A_f(x_1,x_2,y)$
de~$\Lang{L_A}$ telle que $\Th{Q}\Proves
!A_f(\num{n_1},\num{n_2},\num{m})$ dès que $f(n_1,n_2)=m$.
Ainsi, $\Th{Q}$ démontre que la valeur de~$f$ en $n_1,n_2$
est~$m$. Nous exigerons en réalité davantage : que la théorie
démontre qu'aucun autre nombre n'est cette valeur.\footnote{Il s'agit ici de fonctions totales. La représentation exige ces preuves pour chaque tuple fixé de numéraux ; elle n'exige pas une preuve unique, dans la théorie, de la totalité de la fonction pour tous les arguments.} Il en va de
même des relations !!{decidable}s. Les deux définitions suivantes
précisent ces exigences.

\begin{defn}
Une !!{formula}~$!A(x_1,\dots,x_k,y)$ \emph{!!{represents}} la
fonction $f\colon\Nat^k\to\Nat$ dans~$\Gamma$ si et seulement si,
lorsque $f(n_1,\dots,n_k)=m$, on a :
\begin{enumerate}
\item $\Gamma \Proves !A(\num{n_1}, \dots, \num{n_k}, \num{m})$, et
\item $\Gamma \Proves \lforall[y](!A(\num{n_1}, \dots, \num{n_k},
y) \lif y = \num{m})$.
\end{enumerate}
\end{defn}

\begin{defn}
Une !!{formula}~$!A(x_1,\dots,x_k)$ \emph{!!{represents}} la
relation $R\subseteq\Nat^k$ \emph{dans~$\Gamma$} si et seulement si\footnote{Précision éditoriale : la source sous-entend ici~$\Gamma$ ; nous rappelons explicitement que la représentation est relative à la théorie.} :
\begin{enumerate}
\item si $R(n_1,\dots,n_k)$, alors $\Gamma \Proves !A(\num{n_1},
\dots, \num{n_k})$;
\item si $R(n_1,\dots,n_k)$ ne vaut pas, alors $\Gamma \Proves \lnot
!A(\num{n_1}, \dots, \num{n_k})$.
\end{enumerate}
\end{defn}

Une théorie est « assez forte » pour les théorèmes d'incomplétude si
elle !!{represents} toutes les fonctions calculables et toutes les
relations !!{decidable}s. $\Th{Q}$ et ses extensions satisfont cette
condition, mais il nous faudra du temps pour le démontrer. C'est un
fait non trivial sur ce que~$\Th{Q}$ peut prouver, d'autant que cette
théorie ne possède qu'un petit nombre d'axiomes. $\Th{Q}$ est
pourtant très faible. Bien qu'il soit difficile d'établir qu'elle
représente toutes les fonctions calculables, la plupart des théories
intéressantes sont plus fortes qu'elle : elles démontrent au moins
tout ce qu'elle démontre. Toute extension de~$\Th{Q}$ représente donc
elle aussi toutes les fonctions calculables et les relations
!!{decidable}s. Beaucoup de théories intéressantes remplissent
ainsi les hypothèses des théorèmes d'incomplétude. En particulier,
toute théorie censée formaliser « toutes les mathématiques » doit
rendre compte des calculs élémentaires ; nous demanderons ici qu'elle
étende~$\Th{Q}$, en plus des autres hypothèses des théorèmes.\footnote{Précision éditoriale : démontrer des calculs sur des numéraux ne suffit pas, à lui seul, à déduire tous les axiomes universels de~$\Th{Q}$. Pour une théorie formulée dans un autre langage, le critère analogue passe par une interprétation de l'arithmétique requise ; les hypothèses des théorèmes doivent alors être vérifiées dans ce cadre.}

\end{document}
